% Generated draft: requires mathematical and source review.
\documentclass[11pt]{article}
\usepackage[margin=25mm]{geometry}
\usepackage{fontspec}
\setmainfont{Latin Modern Roman}
\usepackage{amsmath,amssymb,mathtools,unicode-math}
\setmathfont{Latin Modern Math}
\usepackage{enumitem,needspace,xurl,hyperref}
\setlength{\parindent}{0pt}
\setlength{\parskip}{5pt}
\setlength{\emergencystretch}{3em}
\newcommand{\N}{\mathbb{N}}
\newcommand{\Z}{\mathbb{Z}}
\newcommand{\Q}{\mathbb{Q}}
\newcommand{\R}{\mathbb{R}}
\newcommand{\C}{\mathbb{C}}
\newcommand{\F}{\mathbb{F}}
\newcommand{\A}{\mathbb{A}}
\newcommand{\PP}{\mathbb{P}}
\newcommand{\E}{\mathbb{E}}
\newcommand{\eps}{\varepsilon}
\newcommand{\dd}{\,\mathrm{d}}
\begin{document}
\section*{UnsolvedMath ID: 20000986}
\textbf{Source problem code:} catalogue entry
\subsection{Definitions and mathematical statement}
The original asks how Baker--Norine's graphical Riemann--Roch theorem might extend to cell complexes, possibly through higher-dimensional critical or superstable configurations and higher-dimensional algebraic Riemann--Roch ideas. It does not specify a higher-dimensional divisor rank, degree, effective cone, canonical class, or target formula. The supplied partial result compares a published simplicial multidegree with a coordinate-free lattice quotient; it does not supply a full higher-dimensional Riemann--Roch theorem.

Fix a finite oriented cell complex and degree \(i\geq1\), with integral cellular boundary maps \(\partial_j:C_j\to C_{j-1}\). Give cellular bases their standard integral inner products. Write
\[
 L=L_i^{\mathrm{ud}}=\partial_{i+1}\partial_{i+1}^{\mathsf T},\qquad
 M=\ker_\Z L=\ker_\Z\partial_{i+1}^{\mathsf T},\qquad
 K_i=\ker_\Z\partial_i/\operatorname{im}_\Z L.
\]
Let \(T(K_i)\) denote the torsion subgroup, not the full group when rational homology is nonzero. Pairing with \(M\) defines
\[
 \delta:\operatorname{coker}L\to\operatorname{Hom}(M,\Z),\qquad
 \delta([c])(z)=\langle c,z\rangle.
\]
The aid's exact lattice statement is
\[
 0\longrightarrow T(K_i)\longrightarrow\operatorname{coker}L
 \xrightarrow{\delta}\operatorname{Hom}(M,\Z)\longrightarrow0.
\]
The kernel lattice \(M\) is saturated, which makes the pairing map surjective. If \(H_i(X;\Q)=0\), then \(K_i\) is finite and \(T(K_i)=K_i\). Without that hypothesis the full critical quotient must not replace its torsion subgroup in the sequence.

For the Kim--Perkinson comparison, let \(\mathcal H=\{h_1,\ldots,h_s\}\subset M\) be the Hilbert-basis family in their multidegree construction. Its asserted factorization is
\[
 \deg_{\mathrm{KP}}=j_{\mathcal H}\circ\delta,\qquad
 j_{\mathcal H}(\varphi)=(\varphi(h_1),\ldots,\varphi(h_s)).
\]
For the reported injective coordinate embedding it is necessary that the \(h_j\) span \(M\otimes\Q\), and under that explicit condition it is sufficient. The aid reports this injectivity for its intended construction, so its realizable degree lattice is \(j_{\mathcal H}(\operatorname{Hom}(M,\Z))\), canonically identified with \(\operatorname{Hom}(M,\Z)\), not necessarily the whole coordinate lattice \(\Z^s\). In this interpretation the kernel assertion of their Theorem 13 is the torsion kernel in the displayed exact sequence. The source summary does not define the Hilbert-basis cone, its orientation assumptions, or the spanning verification; these are indispensable reference details rather than a property of an arbitrary unspecified Hilbert basis.

The independent orientation obstruction concerns effectiveness. If an additive submonoid \(\mathcal E\subset C_i\simeq\Z^N\) contains every positive cellular basis vector and is invariant under independently reversing each coordinate orientation, then it contains both signs of every basis vector and hence \(\mathcal E=\Z^N\). Thus a fixed coordinatewise nonnegative effective cone cannot remain unchanged under all independent higher-cell orientation reversals without becoming the entire chain lattice. This does not forbid covariantly changing the cone with the orientation, or other definitions of effectiveness omitting these hypotheses. The coordinate-free degree comparison and this no-go lemma do not construct higher superstable objects or the missing Riemann--Roch rank theory.
\subsection{Short English statement}
The supplied partial result identifies a realizable simplicial multidegree lattice with the integral dual of the up-Laplacian kernel and isolates the torsion critical-group kernel. Its comparison needs the actual Hilbert-basis spanning hypotheses. A separate orientation lemma explains why an unchanged coordinatewise effective cone cannot be invariant under all cell reversals, leaving the full Riemann--Roch extension unresolved.
\subsection{Sources}
\begin{itemize}
\item Local registry: permanent UnsolvedMath ID \(20000986\), exact public code \texttt{catalogue entry}; \texttt{aim-combinatorics-notes.json}, \texttt{notes[110]}.
\item AIM workshop \emph{Generalizations of chip-firing and the critical group}, Problem 21, attributed to many participants: \url{https://aimath.org/pastworkshops/chipfiringproblems.pdf}. This URL is supplied provenance; no PDF access is claimed.
\item Original reference: Baker and Norine [7] for graphical Riemann--Roch, without the bibliography entry.
\item Supplied research summary, ``A coordinate-free form of simplicial multidegree and an orientation obstruction,'' attempt 1, cites Kim--Perkinson's published construction and Theorem 13 without a title, date, complete theorem hypotheses, or full citation. Original remarks and literature fields are empty.
\end{itemize}
\end{document}
